\section{Proof-assistant verification of the mathematical contract}
\label{sec:formalization}
The certificate interface has a standalone formalization in Lean~4 \citep{demoura2021-lean4}. Its coverage inventory lists 49 mathematical obligations from Sections~\ref{sec:model}--\ref{sec:implementation}: propagation, supporting cuts, curvature rules, discrete inference, master identity, bound transfer, and the stated mathematical examples. A claim-to-theorem map gives the precise hypotheses and declarations for each obligation. The formalization verifies the mathematical contract and includes executable Lean checkers for typed discrete proofs and rational quadratic matrices. It does not establish that the Python implementation refines those checkers, and no benchmark certificate is checked by Lean.

\subsection{Nonlinear cuts and recognition rules}
The coordinate type distinguishes bounded intervals, both half-lines, and the real line, so infinite endpoints never enter real arithmetic. The proofs establish the exact suprema in Table~\ref{tab:shift}, including attainment and the necessary and sufficient finiteness conditions, as well as the enclosure tests in Table~\ref{tab:enclosure}. Rational propagation instructions preserve every mixed-integer feasible point through any finite admitted sequence, including integer rounding, binary intersections, and fixed substitution. Convexity and a justified derivative along each feasible segment imply the supporting inequality. The cut theorem combines that support with genuine value and residual-slope enclosures and the rational intercept test to derive affine underestimation on the whole box. Interval lemmas prove the finite endpoint arithmetic and the sign-dependent rules for unbounded coordinates; they do not verify the execution of an interval library.

The curvature proofs cover the stated composition rules with their domain and sign conditions, quadratic convexity and concavity, PSD homogenization and square-root convexity, real-exponent monomials, and one-variable linear fractions. Quadratic elimination is a structurally recursive rational checker with soundness and completeness proofs for symmetric matrices, including singular and zero-pivot cases. The monomial development proves the actual second derivative in every direction and its displayed Hessian matrix expression, the three exponent criteria, the supporting block identities, and the continuous boundary extensions. These results concern the mathematical functions and checked domains; recognizing their syntax in a loaded expression is a separate implementation obligation.

\subsection{Discrete checking and end-to-end transfer}
The Lean discrete checker consumes rational rows with equality or either weak inequality, declared integer coordinates, rational solution candidates, and a finite list of typed inference steps. It checks row domination, multiplier signs, integral rounding, and exhaustive integer splits. References address only the checked prefix, assumption sets retain their exact indices, and unsplitting removes each branch's own assumption without erasing cross-branch dependencies. The theorem \texttt{Discrete.check\_sound} proves the complete inference invariant from these tests. Every supplied row is checked, including rows after an earlier proving row. Rational solution validation and a proved finite-minimum selection supply the weak incumbent cutoff; \texttt{Discrete.checked\_bound} then proves the unrestricted master bound. Maximization follows by objective negation, and the separate infeasibility checker uses no solution cutoff.

The end-to-end development constructs the feasible extension rather than assuming its inclusion in the master. It connects admitted propagation steps, accepted supporting-cut data, original affine rows, and integrality to the graph point with its free epigraph coordinate and fixed-one constant coordinate. The affine and epigraph objectives preserve the normalized original objective exactly. A typed master-identity check proves equivalence under variable renaming, integrality correspondence, positive rational row scaling, and complete row permutations, together with objective equality. The resulting bound theorems compose those results with the discrete checker; the sign and infeasibility corollaries have the interpretations stated in Section~\ref{sec:bound-transfer}.

Primal completion uses an actual feasible nonlinear witness and an objective enclosure. The proofs establish nonemptiness and boundedness before applying the real infimum, derive the bound chain and gap, and prove attainment when the bounds match. Separate extended-real results cover the empty feasible set. The formal examples also check the coefficient-aggregation arithmetic, the complete slope-rounding calculation, boundary-support cautions, invalid inference examples, and domain counterexamples. Kernel-evaluated discrete examples exercise accepted splitting and rounding, rejected forward references and later invalid rows, retained cross-branch dependencies, and rejected infeasible or fractional incumbents.

\subsection{Validation and trust boundary}
The package pins Lean~4.33.1 and mathlib~4.33.1 \citep{mathlib2020}, with exact transitive revisions in its manifest. The current development was checked with targeted module builds that treat warnings as failures; we report no result of the project-wide verification script. That script checks complete module imports, builds the project, audits every owned declaration for transitive axioms, and replays compiled declarations through the installed Lean kernel. Its axiom allowlist is \texttt{propext}, \texttt{Classical.choice}, and \texttt{Quot.sound}; the replay is another pass through the same kernel, with imported mathlib as the dependency base. An earlier replay of three modules is kept as a historical record and does not cover the current development.

The remaining trust boundary is operational. Python model loading, expression extraction, symbolic differentiation, actual interval enclosures and transcendental evaluation, serialization, VIPR text parsing, mutable row storage, lifetime accounting, and correspondence between Python execution and the typed Lean checkers are not formally verified. Nor do the proofs verify saved benchmark bundles, external-checker behavior, performance measurements, or novelty claims. The recognition rules are sufficient, not complete for all convex expressions or all valid cuts. Accepted nonlinear evidence must still supply genuine derivative and enclosure data for the same interpreted model.
